\section{Contribution to the Dissertation}
\label{sec:checkpoints-contribution}

\subsection{Contribution to Prediction vs.~Coordination}
\label{sec:checkpoints-contribution-prediction-coordination}

The dissertation's central argument distinguishes between prediction-oriented architectures, which
anticipate future states by extending patterns from historical data, and coordination-oriented
architectures, which enforce causal constraints across distributed processes. The checkpoint ledger
is a coordination artifact: it does not predict whether research will be complete; it enforces the
constraint that downstream manufacturing cannot begin until the evidence base satisfies declared
minimum thresholds.

This distinction is most visible in the authorization lists that each ALIVE verdict carries.
PROCESS\_INTELLIGENCE\_ALIVE\_001 authorizes board-level claim delivery and M\&A diligence package
assembly. It does not predict that the M\&A deck will be high-quality; it certifies that the
evidence base is sufficient to make board-admissible claims. The difference is between a
probabilistic forecast (the deck will probably be good) and a causal constraint (the deck may not
be manufactured until the receipt registry contains the required entries). Coordination does not
merely describe the future; it restricts it.

The no-forced-ALIVE rule is the sharpest expression of this contribution. A predictive system can
always produce an optimistic forecast. A coordination system can refuse to issue an authorization
until the evidence exists. The checkpoint ledger implements the latter by making Gate 12
(No Forced ALIVE Verdicts) a first-class audit criterion that must pass for any ALIVE verdict to be
sealed.

\subsection{Contribution to Process Evidence}
\label{sec:checkpoints-contribution-process-evidence}

The checkpoint ledger contributes to the process evidence subsystem of the thesis by demonstrating
that the \texttt{Evidence<T, State, W>} type-law carrier pattern is not merely a programming
construct but a research governance mechanism. Every receipt in the RECEIPT\_REGISTRY.md is an
instance of the Evidence pattern applied to a research artifact: the artifact type is the
manufactured object (paper inventory, lifecycle phase definition, M\&A claim document), the state is
the gate-pass status, and the witness is the named law or standard that certifies the artifact's
admissibility.

The Petri net snapshot in PETRI\_NET\_SNAPSHOT.md provides formal process evidence at the program
level. The decoded final marking---[Decommissioning]:1, all intermediate places at zero---satisfies
the proper completion criterion for a sound Petri net. This is not a claim; it is an evidence
object. The reachability of the [Decommissioning] state from the initial [Design] state, with no
deadlocks, confirms that the research lifecycle process model is sound in the Van der Aalst sense:
every case that enters the process can complete, and no case is permanently stuck.

The adversarial evidence in PROCESS\_INTELLIGENCE\_ADVERSARIAL\_V30.1.1\_OMEGA contributes negative
test coverage: the 20-agent swarm evaluation tests whether the evidence base can withstand hostile
attack vectors, recursive logic bombs, and ontological contradictions. All attack vectors were
neutralized, which is process evidence that the research program's boundary enforcement is effective
under adversarial conditions.

\subsection{Contribution to Manufacturing Architecture}
\label{sec:checkpoints-contribution-manufacturing}

The checkpoint system's contribution to manufacturing architecture is the \emph{authorization gate}
pattern: no manufacturing operator runs until its prerequisites are certified by a sealed checkpoint.
This pattern appears throughout the SUBSTRATE\_COMPLETE\_001 document, which records that all six
manufacturing phases (wasm4pm-compat validation, ggen template embedding, rendering engine
activation, wasm4pm module rendering, M\&A deck manufacturing, Blue River Dam integration) were
completed in declared sequence, each phase dependent on the completion of the prior phase.

The receipt-sealed generation records in SUBSTRATE\_COMPLETE\_001 demonstrate that the manufacturing
architecture is auditable end-to-end. Each of the five wasm4pm module receipts (mining, conformance,
replay, lifecycle, Blue River Dam) records the input governance sources, the rendering decisions, the
test results, and the witness markers injected into the rendered code. An auditor can trace any
line in the manufactured code back to the governance document that authorized it, and can verify
the receipt independently by re-running the rendering engine against the same inputs.

The GGEN\_ECOSYSTEM\_INTEL\_ALIVE\_001 checkpoint contributes the boundary audit pattern: four
automated audits (No DTO Flattening, No Tool Smuggling, Feature Isolation, Graduation Boundary)
that enforce the structural separation between the wasm4pm execution engine and the wasm4pm-compat
FFI layer. Audit 1 is currently failing, which demonstrates that the boundary audit pattern can
detect violations rather than merely permitting them to pass unnoticed.

\subsection{Contribution to Capital-Flow Infrastructure}
\label{sec:checkpoints-contribution-capital-flow}

The checkpoint ledger contributes to the capital-flow and settlement infrastructure of the thesis
by establishing the receipt-to-board-claim traceability chain. Every M\&A claim in the \texttt{ma/}
directory is traceable to a receipt in the RECEIPT\_REGISTRY.md. The SLIDE\_TO\_RECEIPT\_MAP.md in
the \texttt{ma/} directory maps each board presentation slide to the specific module receipt that
authorizes its claim: process discovery accuracy to the wasm4pm mining generation receipt, conformance
status to the conformance generation receipt, compliance automation to the Blue River Dam generation
receipt.

The five monetization claim sample receipts (EBITDA rework, working capital/accounts receivable,
risk/SLA compliance, risk/compliance, residual standard) are JSON documents carrying BLAKE3 proof
chains. They represent the data model for production receipts that will be manufactured against
live M\&A target data in subsequent phases. These samples demonstrate that the receipt format is
expressive enough to carry financial claim evidence and that the proof chain mechanism can support
board-admissible M\&A diligence packages.

The ALIVE authorization of board-level claim delivery in PROCESS\_INTELLIGENCE\_ALIVE\_001 is the
specific connection between the checkpoint system and the capital-flow infrastructure: a board claim
is only admissible if its underlying receipt is indexed in the RECEIPT\_REGISTRY.md. This requirement
converts the receipt ledger from a research bookkeeping tool into an enforcement mechanism for
capital-market claim integrity.
